\section{El permanente y el predictor de torres}\label{app:torres}
\label{app:rook}
\paragraph{La reducción G1 [D].}
En la comparación de $\operatorname{per}M$ con $\operatorname{per}M_e$ cambian tanto el tamaño
como el escalado bistocástico. Separamos el error $g_T$ en dos cantidades que podemos calcular
en un mismo escalado. El vector de incidencia $a_e$ tiene un uno en la fila y otro en la columna
de $e$.
Sean $\tau_e=p_e(1-p_e)/(1-2p_e)$ y $S_T=\sum_e\tau_e a_ea_e^\top$.
El producto de sus autovalores positivos es $\operatorname{pdet}S_T$.
Con sumas sobre las casillas permitidas, la normalización en ambos extremos de un incremento es
\[
\begin{aligned}
 H_T(p)&=\sum_e[-p_e\log(p_e/w_e)+(1-p_e)\log(1-p_e)],\\
 L_T(M)&=H_T(p)-\tfrac12\sum_e\log\frac{1-2p_e}{(1-p_e)^2}
         -\tfrac12\log\operatorname{pdet} S_T+\tfrac12\log(4\pi k).
\end{aligned}
\]
Escalar las filas por $u_i$ y las columnas por $v_j$ incrementa $H_T$ y el logaritmo del
permanente en las mismas sumas de logaritmos. Los demás términos dependen de $p$, invariante
bajo este escalado. Así $g_T(M)=g_T(P)$. Para una matriz bistocástica $D$ de tamaño $K$, defina
\[
 b_D=D-\mathbb U_K,\qquad
 \mathcal E(D)=\log\frac{K^K\operatorname{per}D}{K!}
             +\tfrac12\log\det(I-b_D^\top b_D).
\]
El logaritmo del permanente se cancela exactamente en $\Delta=g_T(M)-\mathcal E(P)$. Quedan dos
cálculos. Para $\mathcal E(P)$ hay que hallar la primera corrección al determinante principal
del permanente. Para $\Delta$ basta desarrollar los términos explícitos de la normalización.

\paragraph{La estimación estática $(L^*)$ [D].}
El determinante principal en la clase difusa con brecha se obtiene en~\cite{dw}. El desarrollo
de McCullagh~\cite{mcc} identifica correcciones formales y la estimación de Li~\cite{li} trata
perturbaciones de norma de línea sublineal. Aquí necesitamos que, al extraer el coeficiente de
$1/K$, el error sea $O(K^{-2})$ con una misma constante para toda la clase. Para probarlo,
reunimos los menores del núcleo centrado en el polinomio
\[
 \mathcal R_b(s)=\sum_{\ell=0}^K s^\ell
         \sum_{|I|=|V|=\ell}\operatorname{per} b[I,V],\qquad b=b_D .
\]
Expandir $\operatorname{per}(\mathbb U_K+b)$ y usar
$\mathbb E T^{-\ell}=(K-\ell)!/K!$ para $T\sim\Gamma(K+1,1)$ da
$K^K\operatorname{per}D/K!=\mathbb E\,\mathcal R_b(K/T)$.
Para $b=U\Sigma V^\top$ y $Z$ gaussiano complejo estándar, la regla de Wick identifica
$\mathcal R_b(s)$ con $\mathbb E\prod_a(1+w_a)$, donde
$w=\sqrt{s}(U\Sigma^{1/2}Z,V\Sigma^{1/2}\overline Z)$.
Los términos lineales suman cero. Absorber $\exp(-\sum_a w_a^2/2)$ en la densidad da
\[
 \mathcal R_b(s)=h(s)\widetilde{\mathbb E}\prod_a (1+w_a)e^{-w_a+w_a^2/2},\qquad
 h(s)=\det(I-s^2b^\top b)^{-1/2}.
\]
Para $s$ real cerca de uno, $C=\widetilde{\mathbb E}ww^\top=sB_0(I+sB_0)^{-1}$,
donde $B_0$ tiene bloques diagonales nulos y bloques cruzados $b,b^\top$.
Cada fila o columna aporta un factor que empieza en $1+w^3/3-w^4/4$. En el primer orden hay un
vértice de grado cuatro o dos de grado tres. Hay tres emparejamientos en el vértice cuártico.
En el caso cúbico, seis dan tres aristas entre vértices y nueve dan un lazo en cada uno.
La expresión con sumas sobre todos los índices es
\[
 \frac{\mathfrak f(s)}K=-\frac34\sum_a C_{aa}^2+
 \sum_{a,a'}\left(\tfrac12 C_{aa}C_{aa'}C_{a'a'}+\tfrac13 C_{aa'}^3\right).
\]
Tómese $\delta=\gamma/4$. Para $|s-1|\le\delta$ la brecha da entradas $O(K^{-1})$
y sumas absolutas de fila acotadas para el propagador. El producto exige índices distintos.
La corrección a la suma sobre todos los pares es $-\frac56\sum_a C_{aa}^3=O(K^{-2})$,
pues $a$ recorre las $2K$ filas y columnas. Esta corrección se conserva al orden siguiente.

Sea $\mathcal W_b(s)=\mathcal R_b(s)/h(s)$ y sea $f_m(s)$ el coeficiente de $t^m$ en $\mathcal W_b(st)$.
Denótese por $f_m^{(j)}$ su parte de exceso $j$ (aristas menos vértices). Para
$N=\lceil\kappa\sqrt K\rceil$, al restar la clase de exceso uno sumada en todos los órdenes queda
\[
 \mathcal W_b(s)-1-\frac{\mathfrak f(s)}K
 =-\frac56\sum_a C_{aa}^3+\sum_{j=2}^{N-1}\sum_{m\le N}f_m^{(j)}
   +\sum_{m>N}f_m-\sum_{m>N}f_m^{(1)} .
\]
Para $2\le j<N$ el mayorante de cada clase truncada es $2j(j\Lambda/K)^j$, con $\Lambda$
uniforme. Su suma es $O(K^{-2})$, pues el cociente entre mayorantes es menor que $1/2$
hasta ese corte para $K$ grande. La cola del interpolante se acota por $C e^{C_1\sqrt K}r^{-N}$
con $r>1$. Para la de exceso uno, los diagramas tienen dos o tres aristas. Distribuir el orden $m$
entre $d$ aristas da $\binom{m-1}{d-1}$ términos con decaimiento $\theta^{m-d}$, donde
$\theta=(1+\delta)\max(1-\gamma,1/2)<1$. Así la cola está acotada por
$C K^{-1}(N+1)^2\theta^{N-2}$. Con $\kappa\log r>C_1$, ambas colas decrecen exponencialmente
en $\sqrt K$, uniformemente para $s\in[1-\delta,1+\delta]$ y para toda la clase.
En el exterior de $|K/T-1|\le\delta$ se usa
$|\mathcal R_b(s)|\le\mathbb E(1+s\sum_l\sigma_l|Z_l|^2/K)^K$.
La densidad gamma absorbe la potencia por traslación. Sus colas y
$\sum_l\sigma_l\le\sqrt{K(M_0-1)}$ dan $O(\exp(-cK+C\sqrt K))$.
Como $\mathbb E(K/T)=1$ y $\mathbb E(K/T-1)^2=1/(K-1)$, la expansión con
momentos tercero y cuarto de orden $O(K^{-2})$ da
\[
 \mathcal E(D)=\mathfrak f_{\rm e}(KD)/K+O(K^{-2}),\qquad
 \mathfrak f_{\rm e}=\mathfrak f(1)+h''(1)/(2h(1)).
\]
El término con $h''(1)$ recoge la variación del determinante al promediar en $K/T$. Para obtener
el coeficiente del error $g_T$ falta sumar a esta expresión la contribución de $\Delta$.
\subsection{El coeficiente del perfil y su estabilidad}
Para sumar las dos correcciones de G1 conviene expresarlas en el mismo perfil $\nu=kP$.
Al componer núcleos usamos las medidas uniformes de filas y columnas, de modo que cada suma
sobre una fibra queda dividida por el tamaño del tablero. Escriba $B=\nu-1$, $A_R=BB^\top$, $A_C=B^\top B$,
$T_X=A_X(I-A_X)^{-1}$, $U=T_RB=BT_C$ y $V=B+U$. Para $X=R,C$ sea
$j_X(t)=(T_X-A_X)(t,t)$, $t_X(t)=T_X(t,t)$ y $j_0=\int j_R=\int j_C$.
Los cubos en integrales dobles son potencias de los valores de los núcleos. En esta notación,
el funcional que entra en (T) es
\begin{align*}
 \mathfrak F={}&\tfrac12j_0^2+\operatorname{tr}(T_R^2)-\tfrac32 j_0
    -\tfrac34\sum_X\int j_X^2+\langle j_R,V j_C\rangle\\
 &-\tfrac12\sum_X\langle j_X,T_Xj_X\rangle-\tfrac13\sum_X\iint T_X^3
    +2\iint BU^2+\tfrac23\iint U^3,\qquad
 \Phi=\mathfrak F-\tfrac1{12}.
\end{align*}
Sea $\bar e$ la proyección de $\nu$ en $L^2(\nu)$ sobre el ortogonal de las funciones
aditivas de fila y columna. Ponga $m_R=\int\nu^2\,dy$, $m_C=\int\nu^2\,dx$. La parte
determinista del coeficiente es
\begin{align*}
 \delta_1={}&\tfrac1{12}+\tfrac56\iint\nu^3+\tfrac12\iint\nu\bar e^2
 -\iint V\nu\bar e-\tfrac12\iint\nu^2\\
 &+\tfrac12\iint\nu^2(t_R+t_C)-\iint\nu^2V-\tfrac14\sum_X\int m_X^2 .
\end{align*}
Para obtener $\delta_1$ de la normalización, relacionamos el punto estacionario $p$ con el
escalado $P$. Existen potenciales $\alpha,\beta$ con
$P_{ij}=\exp(\alpha_i+\beta_j)p_{ij}(1-p_{ij})$. Multiplicar el logaritmo por $P_{ij}$
y sumar elimina el calibre. Sustituyendo en la expresión exacta de $\Delta$ y expandiendo
la entropía y el pseudodeterminante se obtiene $\Delta=\delta_1/k+O(k^{-2})$.
La expansión emplea la difusión para acotar las sumas de potencias y la brecha para
acotar las inversas. Para verificar el coeficiente, úsense $V B^\top=T_R$, $B^\top V=T_C$,
$T_X A_X=T_X-A_X$ y que $B$ tiene marginales cero. Si $g=\iint B^2$ y
$\gamma_X(t)=A_X(t,t)$, la proyección anterior se escribe
\[
 \bar e=B+g-\rho_R-\rho_C,\quad
 \rho_R=(I+T_R)\gamma_R-V\gamma_C,\quad
 \rho_C=(I+T_C)\gamma_C-V^\top\gamma_R .
\]
Sustituir esta expresión y $t_X=\gamma_X+j_X$ en $\delta_1$ y en la fórmula estática da
$\mathfrak f_{\rm e}+\delta_1=\Phi$ en cada malla. Por tanto
$R=R_{\rm est}+R_{\rm det}$, con $R_{\rm est}=\mathcal E(P)-\mathfrak f_{\rm e}/k$ y
$R_{\rm det}=\Delta-\delta_1/k$. Ambos son $O(k^{-2})$, lo que prueba (T). En el cociente de
dos permanentes aparece la diferencia de sus errores. La descomposición permite estimar esa
diferencia por separado para $R_{\rm est}$ y $R_{\rm det}$.

\paragraph{La estabilidad estática $(L^{*+})$ [D].}
En $(L^*)$ bastaba acotar todos los diagramas de exceso al menos dos. Una vez cambiado el
tablero, hay que reconocer cuáles de sus contribuciones sobreviven en el menor. Consérvense
los diagramas de excesos dos y tres, con la corrección de índices coincidentes
$-\frac56\sum C_{aa}^3$. Los mismos mayorantes dejan una cola $O(k^{-4})$.
El promedio gamma y el logaritmo dan $R_{\rm est}=f_2/k^2+f_3/k^3+O(k^{-4})$, con
$f_2,f_3$ uniformemente acotados. Para pasar a un menor, el defecto de masa por fibra es
$O(k^{-1})$. Los potenciales que lo corrigen tienen norma supremo $O(k^{-1})$.
Al quitar una sola casilla el defecto se concentra en dos fibras. La inversa en el calibre tiene
norma supremo acotada y entradas fuera de la diagonal $O(k^{-1})$. Los potenciales son
$O(k^{-1})$ en las dos fibras y $O(k^{-2})$ en el resto. El argumento de perturbación conserva
la clase con cotas $2M_0,\gamma/2$. Al expandir la diferencia de un producto de propagadores
se marca la arista cambiada. Las sumas restantes se acotan sobre un árbol generador, usando
la norma absoluta de fila y la cota de entrada. La restricción del índice a la fibra afectada
ahorra un factor $k^{-1}$ en el borrado localizado. Los coeficientes $f_2,f_3$ cambian en
$O(k^{-1})$ bajo el menor y en $O(k^{-2})$ bajo el borrado. Al variar también los denominadores,
esto da los órdenes $k^{-3}$ y $k^{-4}$ de los respectivos incrementos de $R_{\rm est}$.

\subsection{Una casilla marcada y el incremento del determinante [D]}
Queda el incremento de $R_{\rm det}$. Se puede medir el efecto de $e=(r,c)$ sobre $L_T$
variando su peso o comparando el tablero con el menor tras ocuparla. Para $e$ permitida, sean
$\pi^{\rm der}_e=\partial_{\log w_e}L_T=p_e(1+\hat c_e)$,
$\varepsilon_e=p_e/P_e-1$ y
$r_e=L_T(M_e)-L_T(M)-\log(\pi^{\rm der}_e/w_e)$.
El Teorema A de torres da $|\hat c_e|\le C/k$ y $|r_e|\le C/k^2$. Así $r_e$ mide hasta
segundo orden la discrepancia entre ambas operaciones. Usando el escalado
del menor $P'=\operatorname{diag}(x')P^{(e)}\operatorname{diag}(y')$, donde $P^{(e)}$ es la
submatriz de $P$, se obtiene la identidad exacta
\begin{align*}
 R_{\rm det}(M_e)-R_{\rm det}(M)&=-r_e-\mathsf E_e,\\
 \mathsf E_e={}&\log(1+\varepsilon_e)+\log(1+\hat c_e)
          +\sum_{i\ne r}\log x'_i+\sum_{l\ne c}\log y'_l\\
 &+(k-1)\log(1-1/k)
  +\tfrac12\log\frac{\det(I-b_{P'}^\top b_{P'})}{\det(I-b_P^\top b_P)}\\
 &+\frac{\delta_1((k-1)P')}{k-1}-\frac{\delta_1(kP)}k .
\end{align*}
En efecto, $\Delta=-L_T(P)+\log(k!/k^k)-\frac12\log\det(I-b_P^\top b_P)$.
La covariancia de $L_T$ bajo escalado da las dos sumas de potenciales. El signo del término de
tamaño sale del cociente de los dos factoriales. Para demostrar $(A^+)$ debemos ganar una potencia
de $k$ en la suma firmada $r_e+\mathsf E_e$. La expresión exacta de $L_T$ conserva el
desplazamiento cuadrático del punto estacionario, habitualmente denotado por $\zeta$. Acotar los
once restos de segundo orden en valor absoluto perdería su cancelación.

\paragraph{Parámetros independientes.}
Para conservar la suma firmada, situamos los dos tableros en una misma familia. Variamos por
separado la masa de la fila $r$, la de la columna $c$, el tamaño inverso y el peso de la casilla
$e$. Fije $W=kP$ y tome cada fila y columna como un tipo. Con $u=1/k$, sean $A_-,B_-$ las
probabilidades uniformes que omiten respectivamente la fila y la columna distinguidas. Ponga
\[
  A_s=(1-s)A_-+s\delta_r,\qquad B_t=(1-t)B_-+t\delta_c,\qquad h=1/N .
\]
Multiplique sólo $W_{rc}$ por $v$. En estas medidas los núcleos escalado y estacionario resuelven
$N=W^v\exp(a_i+b_l)$ y $q(1-hq)=W^v\exp(z_i+y_l)$ con marginales uno.
La inversa de la derivada, fijado el calibre, está uniformemente acotada en norma supremo.
La norma cuadrática se controla por la brecha, y la parte integral transforma norma cuadrática
en supremo por $W\le M_0$. Perturbar las medidas cuesta $O(|s-u|+|t-u|)$, y el peso
marcado cuesta $O(\max(|s|,|t|))$ para $|v|<2$. La iteración con la inversa del centro
converge en discos fijos de clase para $h,s,t$ y da la rama holomorfa.
Los determinantes de los núcleos centrados $V$ que intervienen admiten la forma
\[
 D(V)=\sum_{\ell\ge0}\frac{(-1)^\ell}{(\ell!)^2}
  \int\det[V(i_a,l_b)]_{a,b\le\ell}^{\,2}\,
               dA_s^{\otimes\ell}dB_t^{\otimes\ell}.
\]
Cauchy--Binet y Hadamard dan un mayorante $\ell^\ell C^\ell/(\ell!)^2$. La brecha y la
norma de Hilbert--Schmidt del centro dan un suelo positivo para el determinante allí.
La serie y Cauchy permiten tomar su logaritmo en discos uniformemente menores. En la expresión
de $\Delta$ se sustituye $\log(N!/N^N)-\frac12\log(4\pi N)$ por
$-1/h-\frac12\log2+h/12$ y se guarda aparte el resto $\eta_N=O(N^{-3})$ de Stirling.
La expresión resultante $D_*$ tiene $h^2D_*$ acotada y holomorfa en esos discos.
Para masas racionales positivas, replicar cada tipo $N A_s(i)$ o $N B_t(l)$ veces
preserva el operador de tipos y añade sólo valores singulares cero. Aplicando la parte
determinista de G1 a esas réplicas queda $D_*=\delta_1/N+O(N^{-2})$. Así los tres primeros
coeficientes de $h^2D_*$ se anulan y el de grado tres es $\delta_1$. Por continuidad e identidad
holomorfa, estos coeficientes coinciden para todas las masas. Si $|h|\le r/2$ en un disco
de radio fijo $r$ donde $|h^2D_*|\le M$, el resto tras el grado tres, dividido por $h^4$,
está acotado por $\sum_{j\ge4}Mr^{-j}(r/2)^{j-4}=2M/r^4$. Resulta
$R_{\rm det}=\eta_N+h^2D_2(h,s,t,v)$, con las derivadas primeras de $D_2$ acotadas por
la misma fórmula integral en discos interiores.
Sobre una cara de masa cero se replica sólo la parte visible.
Su calibre se elimina en las sumas de potenciales y se identifica su propio $\delta_1$.

El paso al menor sustituye $(h,s,t)=(u,u,u)$ por $(u/(1-u),0,0)$. La variación de $h$
es $O(u^2)$ y la de cada masa es $O(u)$, de donde $\Delta_e R_{\rm det}=O(u^3)$.
La identidad de arriba prueba $(A^+)$ con el signo indicado. En la cara $s=0$ o $t=0$ las
expresiones integradas son independientes de $v$. Por ello el desarrollo de $\partial_vD_2$
carece de los términos con potencia cero de $s$ o $t$. Dividir su serie por $st$ da una
función $H$ acotada en discos interiores por la serie geométrica de las cotas de Cauchy.
En $s=t=u$ la derivada de $R_{\rm det}$ respecto a $v$ es $O(u^4)$ hasta $v=0$.
Integrar sobre $[0,1]$ da el borrado de $(D^+)$. Al sumar los incrementos de $(L^{*+})$ se
obtiene $(T^+)$ para el resto total $R$. Así puede compararse la aproximación de un mismo par
antes y después de pasar a un menor adicional que lo conserve.

\subsection{De la marginal doble al cuadrado}
\label{app:pair}
Para reintroducir una restricción diagonal necesitamos la probabilidad de que las dos torres
ocupen sus extremos. La estimamos primero con el cociente de predictores $\pi^T$ y luego con
la expresión local $\pi_1$ en el tablero de partida. Sea $\Pi$ el proyector ortogonal sobre el
modo fila menos columna. Para el punto de torres ponga
$G^0=(S_T+\Pi)^{-1}$, $\mathsf G(i,j)=a_i^\top G^0a_j$. La clase estacionaria requiere
$\max p_e\le C/k$ y $\lambda_{\min}(S_T+\Pi)\ge\gamma_A>0$. La clase de escalado anterior
implica estos márgenes para $k$ grande. Para un par $E=\{e,f\}$ en filas y columnas distintas
defina
\[
 \pi^T=w_ew_f e^{L_T(M_{ef})-L_T(M)},\quad
 \pi_1=p_ep_f\{1-\mathsf G(e,f)+\hat c_e+\hat c_f\},\quad
 X=p_ep_f(p_e+p_f)\mathsf G(e,f).
\]
En la máscara binaria el cociente $\pi^{\rm ex}/\pi^T$ es exactamente
$\exp\{g_T(M_{ef})-g_T(M)\}$. Retirar dos fibras de cada familia altera sus sumas en
$O(k^{-1})$. El sistema dual en el calibre fijo recupera las masas con potenciales $O(k^{-1})$.
La brecha se reduce a lo sumo en $O(k^{-1})$. Por Schur e identidad de resolvente las entradas
retenidas de las matrices de Green cambian en $O(k^{-2})$. Suprimir los índices en las
contracciones normalizadas da $|\mathfrak F(M_{ef})-\mathfrak F(M)|=O(k^{-1})$, evaluando
$\mathfrak F$ en el perfil de cada máscara. En consecuencia
$\Delta_E g_T=g_T(M_{ef})-g_T(M)=O(k^{-2})$ por (T). Aplicando dos veces la inserción de
torres y la estabilidad de $p,\hat c$ resulta $\pi^T=\pi_1+O(k^{-4})$ (Corolario C de torres).
Esto da $\pi^{\rm ex}-\pi_1=O(k^{-4})$, la afirmación (R0).

\paragraph{La tercera casilla [D].}
Al ocupar una tercera casilla $c$ en otra fila y otra columna, el par sigue disponible en el
menor binario $M_c$. Aplicar (R0) a ambas máscaras sólo daría $O(k^{-4})$ para la diferencia.
Compare $M,M_{ef},M_c,(M_c)_{ef}$ en $\Delta_c\Delta_E(\Phi/k)$ mediante dos parámetros
de masa. La eliminación de $E$ tiene parámetro
$s$, la de $c$ tiene parámetro $t$ y el tamaño inverso es $k^{-1}/(1-2s-t)$.
Los núcleos e inversas que definen $\Phi$ están acotados en discos complejos fijos. El cociente
de la expresión por $1/k$ tiene derivada mixta acotada allí. Integrar en $[0,1/k]^2$
da $|\Delta_c\Delta_E(\Phi/k)|=O(k^{-3})$. La misma cota para el resto sigue de dos
aplicaciones de $(T^+)$. Entonces, en
\[
 \pi^{\rm ex}-\pi_1=(\pi^T-\pi_1)+\pi^T(\exp(\Delta_E g_T)-1)
\]
el segundo sumando cambia en $O(p_ep_f/k^3)$ por la regla del producto y
$\Delta_c\pi^T=O(p_ep_f/k)$. Para el primero se conserva $p_ep_f$ en el producto de las
inserciones. Las segundas diferencias de $\mathsf G,\hat c$ y el cambio $O(k^{-3})$ de $r_e$
en un menor lejano dan $|\Delta_c(\pi^T-\pi_1)|\le Cp_ep_f/k^3$. Esto prueba (R1) con el peso
$\varpi(E)=k^2p_ep_f$. Ahora promediamos sobre la casilla que se ocupa en una fila $r$ ajena a
las del par. La conmutación de tres menores y $\sum_{c\in r}\pi^{\rm der}_c=1$ dan el defecto
de consistencia de $\pi_1$ bajo este promedio,
\[
  D_r=\sum_{\substack{c\in r\\ M_c\supset E}}\pi^{\rm der}_c(M)\pi_1(M_c)-\pi_1(M),
       \qquad |D_r|\le C\varpi(E)/k^5.
\]
Para probarlo, pongamos $R_\pi=\pi^T-\pi_1$. La conmutación da tres sumandos de $D_r$.
El primero es $\pi^T(M)\sum_c\pi_c^{\rm der}(M_{ef})(e^{\Delta_E r_c}-1)$.
La estabilidad (S) del error de inserción da $\Delta_E r_c=O(k^{-3})$.
El segundo es $-\sum_c\pi_c^{\rm der}(M)\Delta_cR_\pi$ sobre las casillas que conservan el par.
Aquí se aplica la cota de incremento anterior. En el tercero,
$R_\pi(M)\sum_{c\text{ perdido}}\pi_c^{\rm der}(M)$, hay a lo sumo dos casillas,
de masa $O(1/k)$ cada uno. Al reunir las tres cotas se obtiene (R2).

En máscaras binarias la
brecha y la factorización del punto estacionario dan un suelo de clase a $kp_e$ en el soporte
para convertir las cotas absolutas de un número fijo de pares en cotas ponderadas. Transportar
esa conversión a vetos arbitrarios o repetirla en grupos crecientes requeriría otra estimación.
Dos bancos de bloques con un tercer vértice tienen $\hat c_e=\hat c_f=\mathsf G(e,f)=0$ y
$\pi^{\rm ex}=\pi_1$ exactamente [O]. En esas familias el error se anula para cualquier
constante positiva de la cota.

\paragraph{Un par y un cuadrado [D].}
En el tablero $B$ que recupera la restricción diagonal de una sola línea $E$, el logaritmo
del conteo cambia en $\log(1-\pi^{\rm ex})$. Para restarle el cambio del predictor, sea
$\lambda=L(B)-L_T(M)$. Con $p=p_e,q=p_f$, la primera expansión da
$\lambda=-\pi_1+O(k^{-4})$ (Teorema 1). La segunda (Teorema 2) da
\[
       \lambda=\log(1-\pi_1)+X+O(k^{-5}).
\]
En el punto de prueba, una actualización de rango dos del operador axial incorpora las dos
casillas y un complemento de Schur incorpora la línea. El desplazamiento de Newton tiene
componente $pq+O(k^{-3})$ en las cuatro fibras y $2pq+O(k^{-3})$ en la línea añadida.
La primera variación y la mitad de la segunda de la acción, el factor local y la traza del
logaritmo dejan una función analítica en siete variables de orden $O(k^{-1})$.
El polinomio de grado cuatro de $\lambda+\pi_1$ es
$-p^2q^2/2+pq(p+q)\mathsf G(e,f)$ [O, coeficientes racionales].
Los denominadores en el origen son uno y la cota de Taylor es uniforme. Así, en el error
$g^\sharp=\log(1-\pi^{\rm ex})-\lambda$ permanecen la diferencia entre los dos logaritmos
y $-X$ (Corolario 5).

El promedio sobre pares ya incluye el efecto de cada diagonal por separado. Si están disponibles
las cuatro casillas, queda por medir cuánto añade imponer ambas restricciones a la vez. Sean
$E'$ el otro par y $F_\square=E\cup E'$. La actualización sobre cuatro casillas y el complemento de
Schur de dimensión dos~\cite{hj} dan una función de dieciocho argumentos pequeños.
Al restar las dos expresiones de una sola línea, tomadas en la misma máscara, su polinomio de
grado cuatro es $2\prod_{i\in F_\square}p_i$ [O]. Éste es el exceso del predictor
(Teorema 3), con resto $O(k^{-5})$. Los dos pares son incompatibles y dejan el mismo menor.
Si las completaciones son positivas, el exceso exacto de evitación es
$\log[(1-2\pi^{\rm ex})/(1-\pi^{\rm ex})^2]$.
Su primer término es $-(\pi^{\rm ex})^2=-\prod_{i\in F_\square}p_i+O(k^{-5})$ por (R0)
y por la factorización de $p_i(1-p_i)$. La contribución conjunta al error es, por tanto,
$-3\prod_{i\in F_\square}p_i+O(k^{-5})$ (Corolario 6). Al tomar la media de frontera se cuenta
esta diferencia una vez por cuadrado completo y se obtiene $-S$ en~\eqref{eq:main}.

